\subsubsection{Consolidated Database}
\label{sssec:consdb}

The Consolidated Database (ConsDB) \citep[][]{DMTN-227} is an SQL-compatible database designed to store and manage metadata for Rubin Observatory science and calibration images.
Metadata are recorded on a per-exposure basis and includes information such as the target name, pointing coordinates, observation time, physical filter and band, exposure duration, and environmental conditions (e.g., temperature, humidity, and wind speed).
These key image metadata are also stored in the Butler Registry (Butler section), however the ConsDB stores additional information  including derived metrics from image processing and information from the EFD transformed into science-ready formats.
% from the time dimension to the exposure dimension.

The ConsDB schema is organized into instrument-specific tables, e.g., AuxTel and LSSTCam, facilitating instrument-specific queries.
Within the LSSTCam schema, data is further structured into tables for individual exposures and detectors.
An example query on the DP2 dataset might retrieve all visits within a specified time range in the $r$-band for a given DP2 pointing.

The ConsDB is hosted at the USDF.
%Following the initial release of DP2, a release of the DP1 exposure-specific ConsDB data will be made available through the RSP, and accessible externally via TAP.
The detailed LSSTCam schema can be found at: \url{https://sdm-schemas.lsst.io/cdb_lsstcam.html}